\section{Background}
\label{sec:rs-background}

The Recursive Spawning Protocol draws from four largely independent research traditions: agent lifecycle management in multi-agent systems, accountability and provenance in distributed systems, the fixed-versus-mutable delegation chain debate, and hierarchical task decomposition in AI systems. This section situates RSP within each tradition and establishes the BX3 Framework as the minimal sufficient architectural substrate for RSP's accountability guarantees.

\subsection{Agent Lifecycle Management in Multi-Agent Systems}

The problem of managing agent creation, delegation, modification, and termination at scale has produced a growing body of protocol standards and architectural frameworks. The Agent Lifecycle Protocol (ALP) working group formalized the stages through which a persistent autonomous agent moves — provisioning, activation, task execution, suspension, and termination — and identified the critical gap: most lifecycle models treat the delegation event as a point-in-time operation rather than an ongoing structural relationship that must persist and remain verifiable throughout the agent's lifetime~\citep{alp2025}. This distinction is central to RSP, which treats the parent-child relationship not as a transient initialization step but as an immutable structural fact that persists for the duration of the child agent's operational life.

The Human Delegation Provenance Protocol (HDP) extended this view by proposing that every delegation chain from a machine actor to a human principal be represented as a cryptographically signed, append-only ledger — so that any node in the chain can be interrogated at runtime to reconstruct its authorization path~\citep{hdp-draft,helixar2025hdp2,hdp-draft-2}. HDP's append-only ledger model is the direct conceptual ancestor of the Fact Layer forensic ledger in RSP. However, HDP does not address the mutability problem: if any actor in the chain can modify the ledger retroactively, the provenance record loses its evidential value. This limitation is resolved structurally by RSP's Fact Layer write protocol, which makes the parent pointer immutable by architectural constraint rather than by policy assertion.

More recent work has formalized the problem of agent delegation in terms of \textit{delegation chains} — ordered sequences of principal-agent relationships from an authorizing human to an executing machine actor. Project Synapse proposed hierarchical agent-centric frameworks for resolving task conflicts in distributed multi-agent systems, emphasizing that delegation chains must maintain semantic coherence across multiple hops of sub-delegation~\citep{synapse}. The Intelligent AI Delegation Framework introduced adaptive coordination mechanisms for multi-agent systems operating in adversarial environments, recognizing that delegation protocols must be robust to principal agent failure, not just subordinate failure~\citep{intelligent-delegation}. RSP builds on this tradition by making the delegation chain the primary architectural primitive rather than an emergent property of policy configuration.

\begin{table}[ht]
\centering
\caption{Comparison of Agent Lifecycle and Delegation Frameworks}
\label{tab:lifecycle-comparison}
\begin{tabular}{p{3.2cm}p{2.8cm}p{2.8cm}p{3.2cm}}
\toprule
\textbf{Framework} & \textbf{Delegation Model} & \textbf{Chain Mutability} & \textbf{Provenance} \\
\midrule
ALP~\citep{alp2025} & Stage-based lifecycle & Mutable by policy & Log-based \\
HDP~\citep{hdp-draft} & Append-only ledger & Mutable by privilege & Cryptographic \\
Synapse~\citep{synapse} & Hierarchical conflict resolution & Mutable at runtime & Semantic coherence \\
Agentic Communities~\citep{agentic-communities} & Role-based community & Mutable by consensus & None specified \\
RSP (this work) & Immutable parent pointer & Immutable by architecture & Fact Layer forensic ledger \\
\bottomrule
\end{tabular}
\end{table}

\subsection{Accountability and Provenance in Distributed Systems}

The problem of maintaining accountable records in distributed systems has deep roots in database theory, blockchain architecture, and AI governance. The W3C Provenance Data Model (PROV-DM) established a formal vocabulary for representing the provenance of entities, activities, and agents — distinguishing between primary实体 (things produced), activities (actions performed), and agents (entities that bear responsibility)~\citep{provenance-standard,moreau2015-prov}. PROV-DM's tripartite ontology of entity-activity-agent maps cleanly onto RSP's node-spawn-event tripartite structure, where each agent node is an entity, each spawn event is an activity, and the parent pointer traces agent-to-agent responsibility relationships.

More recent work has extended provenance models specifically to AI agent interactions. PROV-AGENT proposed a unified provenance model for tracking AI agent interactions in agentic workflows, recognizing that conventional provenance models are insufficient for multi-agent systems where agents spawn subordinate agents, modify their own state, and execute decisions with downstream consequences that must be attributable to a specific authorizing principal~\citep{prov-agent2025}. The key contribution of PROV-AGENT is the recognition that agent interactions are \textit{causally chained} — a spawn event causes the existence of a new agent — and that causal chains require provenance models that capture not just what happened but what caused what to happen. RSP adopts this causal-chain view of agent spawning and enforces it structurally through the immutable parent pointer and forensic ledger.

Immutable audit trails for Non-Human Identities (NHIs) extended this line of work by formalizing the accountability requirements for autonomous AI actors in multi-agent systems — arguing that NHIs require lineage tracking analogous to supply-chain provenance models, with cryptographic commitments at each delegation step~\citep{nhi-lineage}. Certificate Transparency, originally developed for PKI audit logs, established the design pattern of an immutable, append-only log with verifiable checkpoints that can be audited by any third party without requiring trust in the log operator — a model directly applicable to RSP's forensic ledger design~\citep{laurie2013certificate}. The Trustless Provenance Trees framework further developed game-theoretic mechanisms for operator-gated blockchain registries that maintain provenance without requiring central authority, establishing the theoretical grounding for RSP's Fact Layer write protocol that makes parent pointer immutability a structural property rather than a policy assertion~\citep{trustless2026}.

CollectiveOS v4 operationalized these principles in a forensic-grade multi-agent operating system that treats accountability as a first-class architectural concern — demonstrating that provenance can be maintained at runtime in production multi-agent systems, not just in theoretical models~\citep{collectiveos}. The relationship between CollectiveOS and RSP is complementary: CollectiveOS provides the system-level infrastructure, while RSP provides the protocol-level mechanism for enforcing immutable delegation chains within that infrastructure.

\subsection{Fixed vs. Mutable Delegation Chains}

The question of whether delegation chains should be fixed at provisioning or mutable at runtime is one of the central design choices in multi-agent architecture. Mutable delegation chains are the conventional approach: parent nodes can re-delegate to new subordinates, modify scope parameters, and revoke or redirect authority based on operational conditions. This flexibility is argued to be necessary for adaptive multi-agent coordination in dynamic environments~\citep{wooldridge2009multiagent}.

However, mutable delegation chains create a fundamental accountability problem: when the delegation chain can be modified after the fact, there is no architectural guarantee that the chain in effect at any given time corresponds to the chain that was in effect when the agent made its decisions. Post-hoc audit requires reconstructing the chain from logs, and logs are mutable. This creates the \textit{drift problem} that RSP is designed to solve structurally rather than conventionally.

SentinelAgent formalized the security and verifiability requirements for delegation chains in federal multi-agent AI systems, arguing that delegation chains in high-stakes environments must satisfy formal verification conditions — that the chain topology can be proven correct independent of the assertions of any actor within the chain~\citep{sentinelagent}. The formal verification requirement motivates RSP's immutable parent pointer: a pointer that can be modified after provisioning cannot be the subject of formal verification, because the verification assumes a fixed topology that mutable pointers violate.

The fixed-chain position does not require that agent behavior be static — agents can spawn children with refined scope, and chains can grow along authorized paths. What it prohibits is retroactive modification of the chain topology after provisioning, and it prohibits the creation of agent-to-agent relationships that do not trace back to a human anchor. This is the position adopted by RSP. The distinction between mutable and fixed chains is not a question of rigidity versus flexibility — it is a question of where the flexibility is permitted: at spawn time (authorized, scope-refinement-compliant, architecturally bounded) or at any time (retroactive, subject to privilege, structurally unconstrained).

BX3 formalizes this distinction through its Fact Layer: the write protocol of the Fact Layer permits parent pointer establishment at spawn time and prohibits subsequent modification by any actor. This is not a policy constraint enforced by access control lists — it is a structural constraint enforced by the architecture itself, making the immutability of the parent pointer equivalent in strength to the immutability of the blockchain ledger in a permissionless consensus system.

\subsection{Hierarchical Task Decomposition in AI}

Hierarchical task decomposition (HTD) — the process by which a complex task is recursively decomposed into simpler subtasks assigned to specialized actors — is one of the foundational techniques in classical AI planning and remains central to modern LLM-based agent systems. HTN (Hierarchical Task Network) planning formalized the decomposition paradigm, where tasks are refined through sequences of task-reduction rules until primitive actions are reached that can be executed directly by the system~\citep{ghallab2004htn,erpl}. The classical HTN model provides the theoretical grounding for RSP's scope refinement requirement: each spawn event corresponds to a task decomposition step, and the child node's operating scope must be a proper subset of the parent's scope — enforcing that decomposition is always a refinement, never an expansion.

Herbert Simon's bounded rationality theory established the foundational principle that hierarchical decomposition is necessary for tractable decision-making in complex environments — because the computational cost of non-hierarchical planning grows exponentially with problem complexity, while hierarchical decomposition reduces it to polynomial time~\citep{simon}. Simon's insight is directly applicable to multi-agent systems: without hierarchical decomposition, each agent must reason about the full problem space; with hierarchical decomposition, each agent reasons only about its allocated portion of the problem space. RSP's depth bound ($D_{\max}$) is a structural enforcement of this principle — it prevents unbounded decomposition that would eventually produce agents operating outside their authorized scope.

Recent work has extended HTD to LLM-based agent systems. ReAcTree introduced hierarchical LLM agent trees with explicit control flow for long-horizon task planning, demonstrating that recursive task decomposition in LLM agents produces more reliable and interpretable behavior than flat planning approaches — while identifying the critical problem of alignment maintenance across multiple decomposition levels~\citep{reactree2025}. HiPER proposed hierarchical reinforcement learning with explicit credit assignment for LLM agents, formalizing how reward and responsibility propagate backward through a decomposition tree — a framework directly relevant to RSP's forensic ledger, which tracks the same propagation in the opposite direction (downward from principal to agent)~\citep{hiper2025}. StructuredAgent demonstrated that hierarchical planning with structured memory enables long-horizon web tasks that are intractable under flat planning regimes — further evidence that hierarchical decomposition is not just convenient but necessary for complex operational environments~\citep{structuredttt2026}.

FIRMHIVE extended this research specifically to recursive tree-of-agents architectures for autonomous firmware security analysis, demonstrating that recursive spawning of subordinate agents is both necessary and safe when governed by explicit scope constraints and depth limits~\citep{firmhive2025}. The FIRMHIVE result is particularly relevant because it demonstrates that recursive spawning can be implemented safely in adversarial environments — where the failure modes are not mere inefficiency but active security compromise. RSP builds on FIRMHIVE's scope-constraint model while adding the structural immutability guarantee that FIRMHIVE does not provide.

THREAD independently introduced the concept of recursive spawning in the context of LLM reasoning — demonstrating that recursive spawning of thought sub-agents can improve reasoning quality on complex tasks, and providing empirical evidence that spawning depth must be bounded and that termination conditions must be explicit to avoid infinite recursion or reasoning loops~\citep{thread2025}. The THREAD result complements RSP's theoretical guarantees with empirical validation: recursive spawning improves task performance when properly bounded, but unbounded recursive spawning produces reasoning failures that are structurally analogous to autonomous drift in multi-agent systems.

The critical gap that RSP fills relative to this entire line of research is the absence of an \textit{accountability guarantee}. HTD research focuses on decomposition effectiveness — whether the decomposition produces correct or efficient behavior. It does not address the accountability question: whether the decomposition can be traced to a human principal, whether the chain of delegation is immutable, and whether each subordinate agent's operating scope can be verified to be a proper subset of its parent's scope. These are the questions that RSP addresses structurally.

\subsection{The BX3 Framework as Minimal Sufficient Substrate}

The BX3 Framework is a three-layer accountability architecture — Purpose Layer, Bounds Engine, and Fact Layer — designed to ensure that autonomous systems remain structurally accountable to human principals regardless of recursion depth, delegation complexity, or system autonomy level~\citep{bx3framework2026}. RSP is not a separate mechanism layered onto BX3; it is a structural instantiation of the BX3 architecture applied specifically to the problem of recursive agent spawning. Every RSP mechanism maps onto exactly one BX3 layer, and this mapping is not configurable — it is structurally necessary.

The Purpose Layer provides the validation gate at every spawn event. Before a node $n_i$ can spawn $n_{i+1}$, the Purpose Layer must verify that the proposed operating scope $R(n_{i+1})$ is a proper descendant refinement of $R(n_i)$ — strictly narrower, no wider. The Purpose Layer also validates the \textit{spawn necessity declaration}: the originating node must justify why spawning a child is required rather than handling the task directly. These are authorization checks that require semantic reasoning about intent and scope, which is precisely the class of operations that the Purpose Layer is designed to perform in BX3. A system without a Purpose Layer cannot perform these checks structurally — it can only perform them by policy, which is vulnerable to override by actors with sufficient privilege.

The Bounds Engine enforces the depth bound $D_{\max}$, which is a Purpose Layer parameter that defines the maximum number of spawn steps permitted in any chain before escalation to the human anchor is mandatory. The Bounds Engine also enforces the pre-commit hook that validates scope subset relationships at every spawn event before the child node's parent pointer is written to the Fact Layer. These are enforcement operations that must be performed without human intervention at every spawn event, which is precisely the class of operations that the Bounds Engine is designed to perform in BX3.

The Fact Layer establishes and maintains the immutable parent pointer and the forensic ledger. The parent pointer is written exactly once — at provisioning — and the Fact Layer write protocol prohibits any subsequent modification by any actor. The forensic ledger records every spawn event, every scope refinement decision, and every escalation event, providing the complete causal chain from any active node to its human anchor. These are storage and immutability operations that require cryptographic guarantees and append-only semantics, which is precisely what the Fact Layer provides in BX3.

The three-layer mapping is minimal: no layer can perform the operations of another, and no layer is redundant. The Purpose Layer cannot enforce depth bounds because its operations are semantic and context-dependent; the Bounds Engine cannot perform scope refinement validation because its operations are structural and rule-based; the Fact Layer cannot authorize spawns because its operations are evidentiary, not deliberative. This structural necessity — where each capability belongs exclusively to one layer — is the sense in which BX3 is the minimal sufficient substrate for RSP. A system that implements fewer than three layers cannot achieve RSP's guarantees because it lacks the structural capabilities that those layers provide. A system that implements more than three layers may achieve RSP's guarantees but introduces unnecessary complexity; RSP requires exactly what BX3 provides.

The relationship between BX3 and RSP is analogous to the relationship between an operating system kernel and a system call: the kernel provides the structural capabilities; the system call is the specific interface that uses those capabilities to provide a particular guarantee. Just as a system call cannot provide guarantees that the kernel does not support, RSP cannot provide guarantees that BX3 does not structurally enable. This is why RSP is a named pillar of the BX3 Framework rather than an independent mechanism: its guarantees are derived entirely from BX3's three-layer architecture, and those guarantees are unconditionally bound to that architecture.